\section{Recursive children and the finite induction}\label{adp:induction}

We now check that the changed admissions suffice for the entire
recursive argument. All the arithmetic steps in this section are
the exact identities of \cite[Sections 18.3--18.7]{OAI26}.
Their proofs are independent of the lower bound on $\kappa$.
The following list specifies what is imported; the accompanying
argument checks that it is applicable in the enlarged interval.

\begin{center}
\small
\begin{tabular}{@{}>{\raggedright\arraybackslash}p{.26\linewidth}
                  >{\raggedright\arraybackslash}p{.68\linewidth}@{}}
\toprule
Input from \cite{OAI26}&Data retained in the recursive argument\\
\midrule
Section 18.1, (18.4)--(18.10)&
Exact deletion, inducing character, redundant natural zeros,
reflection scales and conductor--support bound.\\
Section 18.3, (18.17)--(18.19)&
Equal products of scales, common mask, one whole-product norm
power, and a support box fixed before the live prime labels.\\
Section 18.4&
Full first Poisson identity and complete common-support allocation;
localization on the genuine sum before full M\"obius inversion.\\
Section 18.5&
One Gauss-row enlargement, all valuation-one, six and seven errors,
and an auxiliary prime pool disjoint from live prime windows.\\
Section 18.6, Lemma 18.2&
Full second correlation, independent child products with their
original prime coefficients and zeros, and exact normalizers.\\
Section 18.7, Lemma 18.3&
Exceptional-row count and masked equal-product lattice
cancellation on the raw signed rectangles.\\
\bottomrule
\end{tabular}
\end{center}

\subsection{The preserved coefficient}

When centering is needed, write $X_1X_2=Y_1Y_2$ and retain the
entire difference of the two profile rectangles. After fixed plain
powers $\mathfrak b_1,\mathfrak b_2$ have been extracted, it is
\[
 D_{\boldsymbol b}(l_1,l_2)=
 \prod_{i=1}^2W_i(q_{\mathfrak b_i l_i}/X_i)
 -\prod_{i=1}^2W_i(q_{\mathfrak b_i l_i}/Y_i).
\]
Within each separated side, the coefficient is the product of this
difference, one common character, one common puncture, one power
$q_{l_1l_2}^{it}$, and the original live prime coefficients.
The first transform adds a Gauss coefficient on the whole product.
The second transform removes it and restores the same prime
coefficients. It does not leave a new coefficient in an individual
prime factor \cite[Lemma 18.2]{OAI26}.

At either transform, common support is extracted with its full
multiplicities, not merely its gcd. A prime supplied by a live prime
factor is frozen if it enters this support, and punctures every
remaining factor on that side. The artificial residual noncoprime
pairs are removed by the \emph{full} finite M\"obius indicator before
forming independent products. Quotients may still contain the
divisor primes; no unproved coprimality is imposed on them.
Both inverse roots and the full Poisson kernel are retained until
their joint Fourier separation.
The coefficient measure is fixed on the whole-product support box
before live labels are chosen, and is common to both rectangles.

Let $N$ be the fixed number of prime factors. Choose $H_N$ to bound
the sum of the absolute logarithms of their relative annular
endpoints, enlarged for all plain windows, reflection multipliers,
fixed arithmetic normalizations, relative dyadic boxes and support
enlargements through the bounded depth.
For every subset of factors, its aggregate support correction is
at most
\[
 \theta_N=H_N/\log Z .
\]
There is one such correction, rather than $N$ separately chosen
length errors. The raw frequency tails are first paid on the genuine
source. Only afterwards is a common retained dyad extended by zero
to its enclosing row ball.

\subsection{Width decrease and child admission}

Here are the relevant exact length identities. Let $c,d$ be the
logarithmic norms of the two first extracted common-support moduli,
and let $R$ be the logarithmic norm of their common primes whose
net character exponents are nonzero modulo six.
Let $E$ be the logarithmic norm of the divisor selected in the
M\"obius expansion of the complementary common row mask. Put
\[
 \widetilde q=q+R+E,\qquad
 K_0=2A-c-d+R+E-m .
\]
This is the displayed support union and nominal first frequency
length of \cite[(18.22)]{OAI26}; let $K$ be its retained dyad length.
In one Gauss norm, $w$ is the length removed from its column and
$w_o\le w$ is the new moving support thereby created. Put
\[
 a_0=A-c-w,\qquad J=d-c+(K_0-K)-2w+w_o .
\]
The row-ball increment is denoted by $g$, and the prime-density
normalizer by $Z^{-\ell}$. Without prime factors,
$\ell=0$ and $g=J_++\sigma$.
With prime factors, the enlarged main norm has
$\ell=\sigma/3$, $g=J_++2\sigma$.
Each extracted valuation error instead has
$\ell=0$, $g=J_++\sigma$ and
$0\le w_o\le w\le7\sigma/3$. Errors are not enlarged again.
The auxiliary pool has exponent $\sigma/3$, whereas all live
lengths are at most $\eta<\sigma/6$; disjointness follows after a
threshold depending only on the fixed windows.

In the second complete-support extraction write the two moduli as
$D_2a,E_2b$, and set
\[
 c_2=\log_Zq_{D_2},\quad d_2=\log_Zq_{E_2},\quad
 b_2=(c_2+d_2)/2,\quad G_c=(D_2,E_2),\quad g_2=\log_Zq_{G_c}.
\]
Let $p_2$ be the length of their common radical. Among equal
non-six-divisible common valuations, let $t_2$ count the primes
where the divided frequency is a unit. Let $V_{\rm id}$ be the
product of the corresponding nonunit primes, and put
$V=\log_Zq_{V_{\rm id}}$.
The complete local correlation, including all zero cases, is
partitioned on these unit and nonunit conditions before its scalar
is bounded. With $\widetilde q$ the full moving support left by the
first transform, the exact child widths are
\begin{align}
 m'&=2a_0-K-g-g_2-V,\qquad
 q'=\widetilde q+w_o+t_2+V,\label{adp:child-widths}\\
 M'&=m'+q'=M+J-g-g_2+t_2\le M-\sigma .
 \label{adp:width-drop}
\end{align}
This is \cite[(18.32)]{OAI26}; $g_2\ge t_2$ and $g-J\ge\sigma$.
The finite-correlation normalizers leave the child the allowance
\begin{equation}\label{adp:child-allowance}
 M'+\Delta_{\rm child},\qquad
 \Delta_{\rm child}=b_2-p_2+w+\mathcal B_c+\ell\ge w+\ell ,
\end{equation}
as in \cite[(18.33)]{OAI26}. The nonnegative first-extraction bonus
is $\mathcal B_c=(3c-5d-R)_+/6$.

The two child characters have the forms
\[
 \tau_1(n)\rho(n)\chi_n(G_cV_{\rm id}h'),\qquad
 \tau_1(n)\rho'(n)\chi_n(-G_cV_{\rm id}h'),\quad \rho,\rho'\in\Theta .
\]
Their inducing-character ratio belongs to $\Theta$; this statement
takes no quotient at a zero. Every moving good prime belongs to the
declared union in \eqref{adp:child-widths}. Six-divisible factors
remain common punctures. The remaining choices at $\mathcal S$
are finite ray choices, and the norm powers only translate Mellin
height. Thus all nonexceptional children still lie in the global
family to which \eqref{adp:prime} applies.

For such children, first take the triangle inequality between
the two rectangles, and then apply the support test to each plain
scale. An empty annulus is discarded; every retained subunit
annulus is rescaled to one. The exact clipped normalizers are those
of \cite[(18.34)--(18.37)]{OAI26}. In particular the reductions
may differ between rectangles. Let $A_{\rm act},z_{\rm act}$ be
the actual lengths after all allocations, clipping and mask deletion,
and $M'_{\rm act}$ their enclosing width. For the two transforms set
\[
 \delta_{{\rm fr},j}=\frac{\xi}{2}+
       \frac{\log(C_dC_{{\rm sec},j})}{\log Z}<\xi
       \qquad(j=1,2).
\]
Here $C_d$ bounds the relative width of a frequency dyad and
$C_{{\rm sec},j}$ bounds the ratio of the nominal frequency scales
within its fixed outer sector. The eventual threshold ensures the
displayed inequalities. Retain
$K_0-K\ge-\delta_{{\rm fr},1}$ and
$0\le M'_{\rm act}-M'\le\delta_{{\rm fr},2}$.
The entire admission defect is
\[
 F_{\rm act}=(A_{\rm act}-M'_{\rm act}+
                    (6\kappa-1)z_{\rm act})_+ .
\]
The exact ledger \cite[(18.38)]{OAI26}, including the clipped
normalizers, gives
\[
 F_{\rm act}\le6(w+\ell)+\delta_{{\rm fr},1}+2\theta_N .
\]
Remove one prefix of whole prime factors until admission holds,
or until none remains. Since every length is at most $\eta$, its
total removed length $d_z$ satisfies
\[
 \kappa d_z\le F_{\rm act}/6+\kappa\eta
       \le\Delta_{\rm child}
          +\delta_{{\rm fr},1}/6+\theta_N/3+\eta .
\]
This uses only $6\kappa-1\ge0$ and $0<\kappa\le1$.
There is one mesh overshoot. If all prime factors disappear, the
remaining plain lengths use the completed unrestricted zero-factor
estimate. After pairing the two children and their exact coefficients,
the extra edge cost is
\begin{equation}\label{adp:edge-cost}
 \delta_{{\rm fr},2}+\delta_{{\rm fr},1}/6+
             \eta+7\theta_N/3 ,
\end{equation}
the bound of \cite[(18.40)]{OAI26}. The already allocated discrete
and analytic costs remain separate.

\subsection{Exceptional rows}

If the inducing character belongs to $\Theta$, keep the raw signed
rectangles. Neither clipping nor the nonprincipal prime bound is
used on them. After fixing the live prime labels, Lemma 18.2 of
\cite{OAI26} supplies one character, one common polynomial-size
mask, and one norm power in both plains and both rectangles.
The lattice estimate of \cite[Lemma 18.3]{OAI26} has main term
\[
 c_{\vartheta,\mathfrak R_*}X^{1+it}I_i(t),
\]
whose coefficient is independent of $X,t$.
Its two product main terms therefore agree at equal products of
scales and cancel before taking absolute values.

The exceptional character count and the first and second support
payments in \cite[(18.41)--(18.44)]{OAI26} are independent of
$\kappa$. In the centered case their remaining deficit is at most
\[
 A-\frac56M-\frac23v-(L-v)_+\le(A-M)_+,
 \qquad v=c+w+\min(c_2,d_2),\quad L=M/4 .
\]
Indeed its maximum as a function of $v\ge0$ is attained at $v=L$
and equals $A-M$. For positive factors, admission implies $A\le M$.
For the padded zero-factor core, $A-M$ is at most its padding.
This also covers exceptional rows admitted when a partition scalar
or enclosing-ball weight is bounded by one after positivity.
No exceptional row is silently supplied to \eqref{adp:prime}.

\subsection{Choice of mesh, height orders and induction losses}

Prove the zero-factor width bands first and then the positive-factor
bands, each in increasing width. Within a band prove the uncentered
range $A\le5M/6$ before its comparison. Reflection gives arbitrary
zero-factor lengths from the padded core; a comparison calls only
the uncentered assertion already proved at that width.

Before specifying $N$, choose a width floor $\rho$, width step
$\sigma$ and padding $\delta_{\rm pad}$, with the latter small
relative to the first two, such that
\[
 T_{\rm term}:=\rho+\delta_{\rm pad}+\rho/6+
                         3\sigma+5\sigma/3<\varepsilon/4 .
\]
These terms pay, respectively, the small-width estimate, its padding,
the terminal prime cost, and the exceptional and diagonal losses.
If the prescribed initial width is at most $M_{\max}$, put
$D=2+\lceil2M_{\max}/\sigma\rceil$.
Let $C_*\ge1$ bound the fixed number of aggregate affine
perturbations in one stage, enlarged to include
\eqref{adp:clipping}. Choose
\[
 \xi<\min\left\{\frac{\delta_{\rm pad}}2,
       \frac{\rho}{100C_*},\frac{\sigma}{4C_*},
       \frac{\varepsilon}{16C_*D}\right\},\qquad
 \eta<\min\left\{\frac{\sigma}{6},
       \frac{\varepsilon}{16C_*D}\right\},\qquad
 \varepsilon_0<\frac{\varepsilon}{16C_*D}.
\]
The affine admission changes by at most
$|\Delta M|+|\Delta A|+5|\Delta z|$, since
$0\le6\kappa-1\le5$. Apply this to aggregate subset errors.
For the fixed slot system, the threshold on $\theta_N$ and endpoint
logarithms, together with the aggregate tolerance, ensures
$e_c\le\rho/100$ and the comparison margins after \eqref{adp:clipping}.
Every strict child then drops width by at least $\sigma/2$;
there is no same-band cycle.

At the width floor the zero-factor core is estimated by absolute
counting, and \eqref{adp:terminal} with \eqref{adp:prime} pays its
positive factors. Above the floor the second-transform diagonal
costs at most $5\sigma/3$ plus the aggregate errors
\cite[(18.31)]{OAI26}. Together with the uncentered exceptional
count, it first completes $A\le5M/6$ in the current band. Only then
are the comparisons in Section~\ref{adp:numerical} invoked.
Exceptional raw scales are never supplied to a clipped moment
assertion. These are precisely the terminal and same-band calls
used in the induction.

For at most $d$ remaining strict calls use the error envelope
\[
 \mathcal E_d=T_{\rm term}+
                   (d+1)C_*(\eta+\xi+\varepsilon_0).
\]
The preceding bounds close backward induction with total error
less than $\varepsilon$. All surviving subsets use the same mesh.
Small divisor losses are subdivided after $N$ is fixed; they are
not charged as a fixed exponent once per prime or once per slot.
Each discrete label mass is extracted once and its subsequent
measure normalized before taking the smooth norm.

Finally fix the internal derivative and height orders backwards
through this finite induction. Lemmas 4.5 and 4.7 of \cite{OAI26}
bound the full weighted Fourier measures and all Euler derivatives
of the radial kernels uniformly in their scale ratio. Increasing
a derivative order increases a fixed seminorm and height degree,
not a power $Z^{J\xi}$. The polynomial factor in \eqref{adp:prime}
is integrated with these weighted measures. Only after these
internal orders are fixed is the external Fourier-tail order
chosen. Reflection, masks and all finite ray choices preserve
the same bounds. This proves Theorem~\ref{adp:plain}, including
its original uniformity and mesh-before-factor-count assertion.
